% LaTeX Template for AI Problem Analysis Output
% AMC12A 2024 Problem 12 - Strategic Analysis

\documentclass[11pt]{article}
\usepackage[utf8]{inputenc}
\usepackage{amsmath, amssymb, amsthm}
\usepackage{geometry}
\usepackage{xcolor}
\usepackage[most]{tcolorbox}
\usepackage{enumitem}
\usepackage{fancyhdr}

% Page setup
\geometry{margin=1in}
\setlength{\headheight}{14.5pt}
\pagestyle{fancy}
\fancyhf{}
\rhead{AMC12/COMC Problem Analysis}
\lhead{2024 AMC 12A Problem 12}
\cfoot{\thepage}

% Color definitions
\definecolor{problemconfig}{RGB}{230, 242, 255}    % Light blue
\definecolor{strategic}{RGB}{255, 230, 242}       % Light pink
\definecolor{tactical}{RGB}{242, 255, 230}        % Light green
\definecolor{techniques}{RGB}{255, 242, 230}      % Light yellow
\definecolor{verification}{RGB}{255, 230, 230}    % Light red
\definecolor{solution}{RGB}{230, 255, 255}        % Light cyan
\definecolor{competition}{RGB}{242, 242, 255}     % Light purple
\definecolor{gold}{RGB}{218, 165, 32}

% Box styles
\newtcolorbox{problemconfigbox}{
    colback=problemconfig,
    colframe=blue,
    title=PROBLEM CONFIGURATION ANALYSIS,
    fonttitle=\bfseries,
    arc=3pt,
    boxrule=1pt,
    breakable
}

\newtcolorbox{strategicbox}{
    colback=strategic,
    colframe=purple,
    title=STRATEGIC FRAMEWORK APPLICATION,
    fonttitle=\bfseries,
    arc=3pt,
    boxrule=1pt,
    breakable
}

\newtcolorbox{tacticalbox}{
    colback=tactical,
    colframe=green,
    title=TACTICAL METHODOLOGY SELECTION,
    fonttitle=\bfseries,
    arc=3pt,
    boxrule=1pt,
    breakable
}

\newtcolorbox{techniquesbox}{
    colback=techniques,
    colframe=orange,
    title=CONVERSION TECHNIQUES APPLICATION,
    fonttitle=\bfseries,
    arc=3pt,
    boxrule=1pt,
    breakable
}

\newtcolorbox{verificationbox}{
    colback=verification,
    colframe=red,
    title=VERIFICATION STRATEGY DESIGN,
    fonttitle=\bfseries,
    arc=3pt,
    boxrule=1pt,
    breakable
}

\newtcolorbox{solutionbox}{
    colback=solution,
    colframe=cyan,
    title=SOLUTION ROADMAP,
    fonttitle=\bfseries,
    arc=3pt,
    boxrule=1pt,
    breakable
}

\newtcolorbox{competitionbox}{
    colback=competition,
    colframe=purple,
    title=COMPETITION STRATEGY INTEGRATION,
    fonttitle=\bfseries,
    arc=3pt,
    boxrule=1pt,
    breakable
}

\newtcolorbox{keyinsightbox}{
    colback=yellow!20,
    colframe=gold,
    title=KEY INSIGHT,
    fonttitle=\bfseries,
    arc=3pt,
    boxrule=2pt,
    leftrule=5pt,
    breakable
}

\newtcolorbox{warningbox}{
    colback=red!10,
    colframe=red!50,
    title=WARNING: Common Pitfall,
    fonttitle=\bfseries,
    arc=3pt,
    boxrule=2pt,
    leftrule=5pt,
    breakable
}

\newtcolorbox{tipbox}{
    colback=green!10,
    colframe=green!50,
    title=PRO TIP,
    fonttitle=\bfseries,
    arc=3pt,
    boxrule=2pt,
    leftrule=5pt,
    breakable
}

\begin{document}

\title{\textbf{AMC12A 2024 Problem 12\\Strategic Analysis}}
\author{Competition Math Framework v2.1}
\date{\today}
\maketitle

\section*{Problem Statement}

The first three terms of a geometric sequence are the integers $a$, $720$, and $b$, where $a < 720 < b$. What is the sum of the digits of the least possible value of $b$?

\vspace{0.5em}
\textbf{Answer Choices:}\\
$\textbf{(A) } 9 \qquad \textbf{(B) } 12 \qquad \textbf{(C) } 16 \qquad \textbf{(D) } 18 \qquad \textbf{(E) } 21$

\begin{problemconfigbox}
\subsection*{A. Problem Classification}
\begin{itemize}[leftmargin=*]
    \item \textbf{Problem Type}: To Find - optimization with constraints
    \item \textbf{Competition Context}: AMC12A (75 min, 25 problems) - Problem 12 of 25
    \item \textbf{Difficulty Band}: Mid-band (P9-16) - requires geometric sequence properties and optimization
    \item \textbf{Mathematical Domain}: Algebra (sequences) with Number Theory (factorization, divisibility)
    \item \textbf{Estimated Time}: 2-3 minutes for students who recognize the optimization strategy; 4-5 minutes otherwise
\end{itemize}

\subsection*{B. Problem Structure Identification}
\begin{itemize}[leftmargin=*]
    \item \textbf{Given Information}: 
    \begin{itemize}
        \item Geometric sequence with three terms: $a$, $720$, $b$
        \item All three terms are integers
        \item Ordering: $a < 720 < b$
        \item $720 = 2^4 \cdot 3^2 \cdot 5$ (prime factorization)
    \end{itemize}
    
    \item \textbf{Constraints}: 
    \begin{itemize}
        \item $a$, $720$, $b$ form a geometric sequence
        \item All terms must be positive integers
        \item $a < 720 < b$ (strictly increasing)
        \item Need to minimize $b$ subject to these constraints
    \end{itemize}
    
    \item \textbf{Objective}: Find the least possible value of $b$, then compute the sum of its digits
    
    \item \textbf{Hidden Assumptions}: 
    \begin{itemize}
        \item Geometric sequence property: $\frac{720}{a} = \frac{b}{720}$, which gives $ab = 720^2$
        \item Common ratio $r = \frac{720}{a} = \frac{b}{720}$
        \item Since $b > 720$, we need $r > 1$
        \item To minimize $b = 720r$, we need to minimize $r > 1$
        \item For $a$ and $b$ to be integers, $r$ must be a rational number whose denominator and numerator divide $720$
        \item The optimal ratio is of the form $\frac{n+1}{n}$ where both $n$ and $n+1$ divide $720$
    \end{itemize}
    
    \item \textbf{Answer Choice Leverage}: 
    \begin{itemize}
        \item Small digit sums (9-21) suggest $b$ is a 3-digit number
        \item Since $b > 720$, likely $b \in [721, 999]$
        \item Choice (E) 21 is largest, suggesting $b$ might have larger digits or more of them
    \end{itemize}
\end{itemize}

\subsection*{C. Strategic Orientation}
\begin{itemize}[leftmargin=*]
    \item \textbf{Hypothesis}: Given geometric sequence $a, 720, b$ with $a < 720 < b$
    \item \textbf{Conclusion}: Find minimum $b$ and its digit sum
    \item \textbf{Notation}: 
    \begin{itemize}
        \item Let $r$ be the common ratio of the geometric sequence
        \item Then $a = \frac{720}{r}$ and $b = 720r$
        \item Constraint: $r > 1$ (since $b > 720$)
        \item Integer constraint: $r = \frac{x}{y}$ where $y \mid 720$ and $x \mid 720$
    \end{itemize}
    
    \item \textbf{Intuitive Approach}: 
    \begin{itemize}
        \item To minimize $b = 720r$ with $r > 1$, we want $r$ as close to 1 as possible
        \item The closest ratio $> 1$ using factors of 720 is $\frac{n+1}{n}$ for large $n$
        \item Need to find largest $n$ such that both $n$ and $n+1$ divide 720
        \item Check consecutive factors of 720
    \end{itemize}
    
    \item \textbf{Cross-Domain Potential}: 
    \begin{itemize}
        \item \textbf{Algebraic}: Use geometric sequence formulas
        \item \textbf{Number Theoretic}: Prime factorization and divisibility
        \item \textbf{Optimization}: Minimize ratio subject to integer constraints
        \item \textbf{Computational}: Test consecutive factors systematically
    \end{itemize}
\end{itemize}
\end{problemconfigbox}

\begin{strategicbox}
\subsection*{A. Psychological Strategies}
\begin{itemize}[leftmargin=*]
    \item \textbf{Mental Toughness}: The optimization aspect requires careful thinking. Don't rush to test random ratios - develop a systematic strategy.
    
    \item \textbf{Creativity Cultivation}: The key insight is that ``minimize $r > 1$'' translates to finding consecutive factors. This requires thinking about the structure of the problem, not just computation.
    
    \item \textbf{Iterative Refinement}: If first attempts don't work, refine the strategy. If $r = \frac{9}{8}$ seems good, check if there's an even closer ratio like $\frac{16}{15}$.
\end{itemize}

\subsection*{B. Investigation Strategies}
\begin{itemize}[leftmargin=*]
    \item \textbf{Startup Orientation}: 
    \begin{itemize}
        \item Recognize geometric sequence property: middle term is geometric mean
        \item $720^2 = ab$ (key equation)
        \item Factor 720: $720 = 2^4 \cdot 3^2 \cdot 5$
    \end{itemize}
    
    \item \textbf{Get Your Hands Dirty}: 
    \begin{itemize}
        \item Try $r = \frac{3}{2}$: $b = 720 \cdot \frac{3}{2} = 1080$, $a = 720 \cdot \frac{2}{3} = 480$. Valid but large.
        \item Try $r = \frac{4}{3}$: $b = 720 \cdot \frac{4}{3} = 960$, $a = 720 \cdot \frac{3}{4} = 540$. Better!
        \item Try $r = \frac{5}{4}$: $b = 720 \cdot \frac{5}{4} = 900$, $a = 720 \cdot \frac{4}{5} = 576$. Even better!
        \item Pattern: larger denominator (closer to 1) gives smaller $b$
    \end{itemize}
    
    \item \textbf{Wishful Thinking}: ``I wish $r$ could be as close to 1 as possible... What if $r = \frac{n+1}{n}$ for large $n$?''
    
    \item \textbf{Penultimate Step}: Before computing $b$, determine the optimal ratio $r$
    \begin{itemize}
        \item Find largest consecutive integers that both divide 720
        \item Check: do 15 and 16 both divide 720?
        \item $720 / 15 = 48$ \checkmark
        \item $720 / 16 = 45$ \checkmark
        \item So $r = \frac{16}{15}$ is optimal
    \end{itemize}
    
    \item \textbf{Make It Easier}: 
    \begin{itemize}
        \item Instead of testing all possible $r = \frac{x}{y}$, focus on consecutive factors
        \item List factors of 720: $1, 2, 3, 4, 5, 6, 8, 9, 10, 12, 15, 16, 18, 20, 24, \ldots$
        \item Look for consecutive pairs: $(1,2), (2,3), (3,4), (4,5), (5,6), (8,9), (9,10), (15,16)$
        \item Largest consecutive pair: $(15, 16)$
    \end{itemize}
    
    \item \textbf{Visualization}: 
    \begin{itemize}
        \item Think of $r$ on a number line: we want the smallest value $> 1$
        \item Ratios like $\frac{16}{15}, \frac{10}{9}, \frac{9}{8}, \ldots$ get farther from 1
        \item The ratio $\frac{16}{15} \approx 1.0667$ is closest to 1
    \end{itemize}
\end{itemize}

\subsection*{C. Argument Construction}
\begin{itemize}[leftmargin=*]
    \item \textbf{Proof Method}: Direct optimization using geometric sequence properties
    
    \item \textbf{Key Lemmas}: 
    \begin{itemize}
        \item \textbf{Geometric Mean Property}: In geometric sequence $a, m, b$, we have $m^2 = ab$
        \item \textbf{Common Ratio}: If $r$ is the common ratio, then $m = ar$ and $b = ar^2 = mr$
        \item \textbf{Minimization}: For $b = 720r$ with $r > 1$ integer-constrained, minimize $r$
        \item \textbf{Consecutive Factors Optimal}: The ratio $\frac{n+1}{n}$ is minimal when $n$ is maximized
        \item \textbf{Divisibility Requirement}: Both $n$ and $n+1$ must divide 720 for integer $a$ and $b$
    \end{itemize}
    
    \item \textbf{Logical Structure}: 
    \begin{enumerate}
        \item Use geometric sequence property: $ab = 720^2$
        \item Express in terms of common ratio: $b = 720r$, $a = 720/r$
        \item To minimize $b$, minimize $r$ subject to $r > 1$ and integer constraints
        \item For integers: $r = \frac{x}{y}$ where $x, y$ divide 720
        \item Optimal form: $r = \frac{n+1}{n}$ (consecutive factors)
        \item Find largest $n$ such that $n \mid 720$ and $(n+1) \mid 720$
        \item Identify $n = 15$: both 15 and 16 divide 720
        \item Compute $b = 720 \cdot \frac{16}{15} = 768$
        \item Sum digits: $7 + 6 + 8 = 21$
    \end{enumerate}
\end{itemize}
\end{strategicbox}

\begin{tacticalbox}
\subsection*{A. Core Mathematical Tactics}
\begin{itemize}[leftmargin=*]
    \item \textbf{Optimization}: Minimize objective function subject to constraints
    
    \item \textbf{Extreme Principle}: Look for boundary or extreme cases (consecutive factors)
    
    \item \textbf{Factorization}: Use prime factorization to find all factors
    
    \item \textbf{Pattern Recognition}: Recognize that $\frac{n+1}{n}$ approaches 1 as $n$ increases
\end{itemize}

\subsection*{B. Domain-Specific Tactics}
\begin{itemize}[leftmargin=*]
    \item \textbf{Algebra (Sequences)}: 
    \begin{itemize}
        \item Geometric sequence formulas
        \item Common ratio properties
        \item Geometric mean relationship
    \end{itemize}
    
    \item \textbf{Number Theory}: 
    \begin{itemize}
        \item Prime factorization: $720 = 2^4 \cdot 3^2 \cdot 5$
        \item Factor enumeration
        \item Divisibility testing
        \item Consecutive factor search
    \end{itemize}
    
    \item \textbf{Optimization}: 
    \begin{itemize}
        \item Constrained optimization
        \item Minimizing ratios
        \item Integer programming concepts
    \end{itemize}
\end{itemize}

\subsection*{C. Crossover Techniques}
\begin{itemize}[leftmargin=*]
    \item \textbf{Algebra + Number Theory}: Use sequence properties with divisibility
    
    \item \textbf{Optimization + Number Theory}: Minimize ratio using factor structure
    
    \item \textbf{Discrete + Continuous}: Minimize continuous function subject to discrete constraints
\end{itemize}
\end{tacticalbox}

\begin{techniquesbox}
\subsection*{A. High-Probability Techniques}
\begin{itemize}[leftmargin=*]
    \item \textbf{Primary Technique - Consecutive Factors (Optimization)}: 
    \begin{itemize}
        \item \textbf{Recognition Cue}: Need to minimize ratio $> 1$ with integer constraints
        \item \textbf{Application}: Find largest consecutive factors of 720: $(15, 16)$
        \item \textbf{Success Rate}: 85-95\% for students who recognize the pattern
    \end{itemize}
    
    \item \textbf{Secondary Technique - Prime Factorization}: 
    \begin{itemize}
        \item \textbf{Recognition Cue}: Need to find factors of 720
        \item \textbf{Application}: $720 = 2^4 \cdot 3^2 \cdot 5$, enumerate all factors
        \item \textbf{Success Rate}: 90\%+ - straightforward computation
    \end{itemize}
    
    \item \textbf{Alternative - Systematic Testing}: 
    \begin{itemize}
        \item \textbf{Recognition Cue}: Can test small ratios if optimization unclear
        \item \textbf{Application}: Try $r = \frac{3}{2}, \frac{4}{3}, \frac{5}{4}, \ldots$ until pattern emerges
        \item \textbf{Success Rate}: 75-85\% - reliable but slower
    \end{itemize}
\end{itemize}

\subsection*{B. Technique Selection Strategy}
\begin{itemize}[leftmargin=*]
    \item \textbf{Optimal Approach}: Factor 720, find largest consecutive factors, compute $b$
    
    \item \textbf{Backup Approach}: Test several small ratios systematically to find pattern
    
    \item \textbf{Quick Check Approach}: Use answer choices to estimate $b$ (digit sum 21 suggests $b$ around 750-800)
    
    \item \textbf{Time Efficiency}: 
    \begin{itemize}
        \item Consecutive factors method: 1.5-2.5 minutes
        \item Systematic testing: 3-4 minutes
        \item Prime factorization + enumeration: 2-3 minutes
    \end{itemize}
\end{itemize}

\subsection*{C. Technique Integration}
\begin{itemize}[leftmargin=*]
    \item \textbf{Integration Plan}:
    \begin{enumerate}
        \item Recognize geometric sequence property (20 seconds)
        \item Set up $ab = 720^2$ and $b = 720r$ (20 seconds)
        \item Factor 720: $2^4 \cdot 3^2 \cdot 5$ (20 seconds)
        \item List factors, find consecutive pairs (30 seconds)
        \item Identify $(15, 16)$ as largest consecutive pair (10 seconds)
        \item Compute $b = 720 \cdot \frac{16}{15} = 768$ (20 seconds)
        \item Sum digits: $7 + 6 + 8 = 21$ (10 seconds)
        \item Verify (20 seconds)
    \end{enumerate}
    
    \item \textbf{Conflict Resolution}: If consecutive factors strategy is unclear, fall back to testing small ratios
\end{itemize}
\end{techniquesbox}

\begin{verificationbox}
\subsection*{A. Verification Level Selection}

Using the 7-level verification system:

\begin{itemize}[leftmargin=*]
    \item \textbf{Selected Level}: Level 3 - Targeted Spot Check (15-30 seconds)
    \item \textbf{Rationale}: 
    \begin{itemize}
        \item Mid-band problem (P12) with clear solution path once strategy is found
        \item Can quickly verify by checking $a$ and $b$ are integers
        \item Can also verify geometric sequence property
        \item 15-30 seconds sufficient for key checks
    \end{itemize}
\end{itemize}

\subsection*{B. Verification Checklist}
\begin{itemize}[leftmargin=*]
    \item \textbf{Divisibility Check}: 
    \begin{itemize}
        \item Does $15 \mid 720$? $720 / 15 = 48$ \checkmark
        \item Does $16 \mid 720$? $720 / 16 = 45$ \checkmark
        \item Both divide evenly
    \end{itemize}
    
    \item \textbf{Integer Check}:
    \begin{itemize}
        \item $b = 720 \cdot \frac{16}{15} = \frac{720 \cdot 16}{15} = \frac{11520}{15} = 768$ \checkmark (integer)
        \item $a = 720 \cdot \frac{15}{16} = \frac{720 \cdot 15}{16} = \frac{10800}{16} = 675$ \checkmark (integer)
    \end{itemize}
    
    \item \textbf{Geometric Sequence Verification}:
    \begin{itemize}
        \item Check: $ab = 720^2$?
        \item $675 \cdot 768 = 518400$
        \item $720^2 = 518400$ \checkmark
        \item Alternatively: $\frac{720}{675} = \frac{720}{675} = \frac{16}{15}$ and $\frac{768}{720} = \frac{768}{720} = \frac{16}{15}$ \checkmark
    \end{itemize}
    
    \item \textbf{Ordering Check}:
    \begin{itemize}
        \item $a = 675 < 720 < 768 = b$ \checkmark
    \end{itemize}
    
    \item \textbf{Minimality Check}:
    \begin{itemize}
        \item Check if there's a larger consecutive pair: Next would be $(16, 17)$
        \item Does $17 \mid 720$? $720 / 17 = 42.35\ldots$ \texttimes
        \item So $(15, 16)$ is indeed the largest consecutive pair
    \end{itemize}
    
    \item \textbf{Digit Sum Check}:
    \begin{itemize}
        \item $b = 768$
        \item Digit sum: $7 + 6 + 8 = 21$ \checkmark
    \end{itemize}
\end{itemize}

\subsection*{C. Confidence Calibration}
\begin{itemize}[leftmargin=*]
    \item \textbf{Target Confidence}: 95\%+
    \item \textbf{Confidence Indicators}:
    \begin{itemize}
        \item[$\checkmark$] Consecutive factors strategy is sound
        \item[$\checkmark$] Both $a$ and $b$ are integers
        \item[$\checkmark$] Geometric sequence property verified
        \item[$\checkmark$] Ordering constraint satisfied
        \item[$\checkmark$] Answer matches choice (E)
    \end{itemize}
    \item \textbf{Red Flags}: None - solution is clean and verifiable
\end{itemize}
\end{verificationbox}

\begin{solutionbox}
\subsection*{A. Step-by-Step Solution}

\textbf{Method 1: Consecutive Factors (Primary)}

\textbf{Step 1: Use geometric sequence property}

For a geometric sequence $a$, $720$, $b$ with common ratio $r$:
\[
\frac{720}{a} = \frac{b}{720} = r
\]

This gives us:
\[
ab = 720^2
\]

Also:
\[
a = \frac{720}{r} \quad \text{and} \quad b = 720r
\]

\textbf{Step 2: Set up optimization}

Since $b > 720$, we need $r > 1$.

To minimize $b = 720r$, we need to minimize $r$ subject to:
\begin{itemize}
    \item $r > 1$
    \item $a = \frac{720}{r}$ is a positive integer
    \item $b = 720r$ is a positive integer
\end{itemize}

\textbf{Step 3: Find prime factorization}

\[
720 = 2^4 \cdot 3^2 \cdot 5
\]

\textbf{Step 4: Determine optimal ratio form}

For $a$ and $b$ to both be integers, $r = \frac{x}{y}$ where both $x$ and $y$ divide 720.

To minimize $r > 1$, we want $r$ as close to 1 as possible. This is achieved when $r = \frac{n+1}{n}$ for the largest possible $n$ such that both $n$ and $n+1$ divide 720.

\textbf{Step 5: List factors of 720}

Factors of 720: $1, 2, 3, 4, 5, 6, 8, 9, 10, 12, 15, 16, 18, 20, 24, 30, 36, 40, 45, 48, 60, 72, 80, 90, 120, 144, 180, 240, 360, 720$

\textbf{Step 6: Find consecutive factors}

Looking for consecutive integers that are both factors:
\begin{itemize}
    \item $(1, 2)$: both divide 720 \checkmark
    \item $(2, 3)$: both divide 720 \checkmark
    \item $(3, 4)$: both divide 720 \checkmark
    \item $(4, 5)$: both divide 720 \checkmark
    \item $(5, 6)$: both divide 720 \checkmark
    \item $(7, 8)$: 7 does not divide 720 \texttimes
    \item $(8, 9)$: both divide 720 \checkmark
    \item $(9, 10)$: both divide 720 \checkmark
    \item $(10, 11)$: 11 does not divide 720 \texttimes
    \item $(15, 16)$: Check: $720/15 = 48$, $720/16 = 45$. Both divide 720 \checkmark
    \item $(16, 17)$: 17 does not divide 720 \texttimes
\end{itemize}

The largest consecutive pair is $(15, 16)$.

\textbf{Step 7: Compute $b$}

With $r = \frac{16}{15}$:
\[
b = 720 \cdot \frac{16}{15} = \frac{11520}{15} = 768
\]

\textbf{Step 8: Verify and find digit sum}

Check: $a = 720 \cdot \frac{15}{16} = \frac{10800}{16} = 675$

Verify geometric sequence:
\[
\frac{720}{675} = \frac{16}{15} \quad \text{and} \quad \frac{768}{720} = \frac{16}{15} \quad \checkmark
\]

Digit sum of $b = 768$:
\[
7 + 6 + 8 = 21
\]

\textbf{Answer: (E) 21}

\subsection*{B. Alternative Approaches}

\textbf{Method 2: Systematic Testing}

Test ratios of the form $\frac{n+1}{n}$ for increasing $n$:

\begin{itemize}
    \item $r = \frac{3}{2}$: $b = 720 \cdot \frac{3}{2} = 1080$, $a = 480$. Check: $480 \cdot 1080 = 518400 = 720^2$ \checkmark
    \item $r = \frac{4}{3}$: $b = 720 \cdot \frac{4}{3} = 960$, $a = 540$. Valid \checkmark
    \item $r = \frac{5}{4}$: $b = 720 \cdot \frac{5}{4} = 900$, $a = 576$. Valid \checkmark
    \item $r = \frac{6}{5}$: $b = 720 \cdot \frac{6}{5} = 864$, $a = 600$. Valid \checkmark
    \item $r = \frac{9}{8}$: Check if both divide 720. Yes! $b = 720 \cdot \frac{9}{8} = 810$, $a = 640$. Valid \checkmark
    \item $r = \frac{10}{9}$: $b = 720 \cdot \frac{10}{9} = 800$, $a = 648$. Valid \checkmark
    \item $r = \frac{16}{15}$: $b = 720 \cdot \frac{16}{15} = 768$, $a = 675$. Valid \checkmark
\end{itemize}

The smallest $b$ is 768 from $r = \frac{16}{15}$.

Digit sum: $7 + 6 + 8 = 21$

\textbf{Method 3: Using $ab = 720^2$}

We have $ab = 720^2 = 518400$.

To minimize $b > 720$, we want to maximize $a < 720$.

Since $ab = 518400$, the largest $a < 720$ that divides 518400 gives the smallest $b$.

Try values near 720:
\begin{itemize}
    \item $a = 719$: $518400 / 719 = 720.86\ldots$ (not integer)
    \item $a = 718$: Check divisibility... (tedious)
    \item ...
    \item $a = 675$: $518400 / 675 = 768$ \checkmark
\end{itemize}

We can verify that 675 is optimal by noting that $675 = 720 \cdot \frac{15}{16}$, and $(15, 16)$ are the largest consecutive factors of 720.

Thus $b = 768$ and digit sum is 21.

\textbf{Method 4: Using Answer Choices}

The digit sums range from 9 to 21. Working backwards:
\begin{itemize}
    \item Digit sum 21: Could be $759, 768, 777, 786, 795, 867, 876, \ldots$
    \item Digit sum 18: Could be $729, 738, 747, 756, 765, 774, 783, 792, \ldots$
\end{itemize}

Test candidates near 720:
\begin{itemize}
    \item $b = 768$: Check if $720^2 / 768 = 675$ is integer. $720^2 = 518400$, $518400 / 768 = 675$ \checkmark
\end{itemize}

Digit sum: $7 + 6 + 8 = 21$

\end{solutionbox}

\begin{competitionbox}
\subsection*{A. Time Management}
\begin{itemize}[leftmargin=*]
    \item \textbf{Estimated Time}: 2-3 minutes total
    \begin{itemize}
        \item Understanding problem and geometric sequence property: 20 seconds
        \item Prime factorization of 720: 20 seconds
        \item Finding consecutive factors: 40 seconds
        \item Computing $b = 768$: 20 seconds
        \item Computing digit sum: 10 seconds
        \item Verification: 20 seconds
    \end{itemize}
    
    \item \textbf{Time Checkpoints}: 
    \begin{itemize}
        \item At 1 min: Should have $ab = 720^2$ and optimization goal
        \item At 1.5 min: Should have found consecutive factors $(15, 16)$
        \item At 2 min: Should have computed $b = 768$
        \item At 2.5 min: Should have final answer
    \end{itemize}
    
    \item \textbf{Priority Level}: HIGH
    \begin{itemize}
        \item P12 is very accessible with right insight
        \item Quick solve if consecutive factors strategy is recognized
        \item Good point value for time invested
    \end{itemize}
\end{itemize}

\subsection*{B. Risk Assessment}
\begin{itemize}[leftmargin=*]
    \item \textbf{Confidence Level}: 95\%+
    \begin{itemize}
        \item Clear optimization strategy
        \item Verification confirms answer
        \item Multiple solution paths available
        \item Answer matches choice (E)
    \end{itemize}
    
    \item \textbf{Abandonment Criteria}: 
    \begin{itemize}
        \item If after 3 minutes no progress $\rightarrow$ test answer choices
        \item If optimization is unclear $\rightarrow$ try systematic testing
        \item Should NOT abandon - problem is accessible
    \end{itemize}
    
    \item \textbf{Flag for Review}: NO (high confidence, verified answer)
\end{itemize}

\subsection*{C. Competition Strategy Notes}
\begin{itemize}[leftmargin=*]
    \item \textbf{Geometric Sequence Recognition}: The phrase ``geometric sequence'' should immediately trigger $\frac{m}{a} = \frac{b}{m}$ and $m^2 = ab$
    
    \item \textbf{Optimization Insight}: ``Least possible value'' combined with ``$r > 1$'' suggests finding ratio closest to 1
    
    \item \textbf{Consecutive Factors Strategy}: 
    \begin{itemize}
        \item This is a standard AMC technique for ratio optimization
        \item Always check for consecutive factor pairs when minimizing ratios
    \end{itemize}
    
    \item \textbf{Answer Choice Strategy}: Digit sum 21 suggests larger digits, pointing toward numbers in 760s-770s range
\end{itemize}
\end{competitionbox}

\begin{keyinsightbox}
\textbf{Key Insight}: To minimize $b = 720r$ subject to $r > 1$ and integer constraints, we need $r$ as close to 1 as possible. The optimal ratio is $\frac{n+1}{n}$ where both $n$ and $n+1$ divide 720. By finding the largest such consecutive pair $(15, 16)$, we get $r = \frac{16}{15}$, yielding $b = 768$ with digit sum 21. This ``consecutive factors'' strategy is a powerful optimization technique for AMC problems.
\end{keyinsightbox}

\begin{warningbox}
\textbf{Warning}: Don't just try random ratios without a strategy! Also, be careful to check that BOTH the numerator and denominator of $r$ divide 720 - otherwise $a$ or $b$ won't be integers. A common mistake is finding $r = \frac{9}{8}$ and stopping, not realizing that $r = \frac{16}{15}$ is closer to 1 and gives a smaller $b$. Always verify that your consecutive pair is the LARGEST possible.
\end{warningbox}

\begin{tipbox}
\textbf{Pro Tip}: For AMC problems involving minimizing ratios with divisibility constraints, the ``consecutive factors'' strategy is incredibly powerful. Always factor the key number (here 720), list factors in order, and look for the largest consecutive pair. This same technique appears in many AMC problems. Also, recognizing that $720 = 16 \times 45 = 16 \times 3 \times 15$ immediately suggests checking if 15 and 16 both divide 720, which they do!
\end{tipbox}

\section*{Final Answer}

\boxed{\textbf{(E) } 21}

\section*{Confidence Assessment}
\begin{itemize}[leftmargin=*]
    \item \textbf{Confidence Level}: 95-98\%
    \begin{itemize}
        \item Geometric sequence property correctly applied: $ab = 720^2$
        \item Optimization strategy is sound: minimize $r > 1$ by finding largest consecutive factors
        \item Verified that 15 and 16 both divide 720
        \item Computed $b = 768$ correctly
        \item Verified geometric sequence property: $675 \cdot 768 = 720^2$
        \item Verified ordering: $675 < 720 < 768$
        \item Checked that $(15, 16)$ is the largest consecutive pair
        \item Digit sum calculation confirmed: $7 + 6 + 8 = 21$
    \end{itemize}
    
    \item \textbf{Verification Applied}: Level 3 - Targeted Spot Check
    \begin{itemize}
        \item Verified divisibility of 15 and 16 into 720
        \item Checked integer values for $a = 675$ and $b = 768$
        \item Confirmed geometric sequence: $\frac{720}{675} = \frac{768}{720} = \frac{16}{15}$
        \item Verified minimality by checking no larger consecutive pair exists
    \end{itemize}
    
    \item \textbf{Flag for Review}: No
    
    \item \textbf{Time Spent}: Approximately 2-3 minutes (optimal for P12 mid-band)
\end{itemize}

\end{document}

